2.12. 设 $t\in S^*$，$S'$ 是 $S$ 的控制 $S_+(t)$ 的双有理变换（2.3）。$t$ 在 $S'$ 上的\underline{全变换}仍记为 $t^*$，它是 $S_+(t)$ 上的除子 $t^*$ 在 $S'$ 上的全变换。

设 $\nu:S^*\longrightarrow\mathbf N$ 是有限支集函数，并假设对每个 $t\in S^*$，

\begin{equation}\tag{2.12.1}
\nu(t)\geq
\sum_{\substack{t'\ \mathrm{proximate}\\\mathrm{to}\ t}}
\nu(t')[k(t'):k(t)]
\end{equation}
（邻接点不等式）。设 $f:S'\longrightarrow S$ 是 $S$ 的双有理变换，对 $\nu$ 的支集中的每个 $t$，均控制 $S_+(t)$；令
\[
\nu^*=\sum_{t\in S^*}\nu(t)t^*
\]
为 $S'$ 上相应的除子。

%%%%%%%%%% source scan idx 29 | folio 22 %%%%%%%%%%

用 $I(\nu)$ 表示将 $\underline{O}_{S'}$ 中定义 $\nu^*$ 的理想沿 $f$ 取正向像所得的 $\underline{O}_S$ 的理想：

\begin{equation*}
I(\nu)=f_*\underline{O}_{S'}(-\nu^*) \quad .
\end{equation*}
由（2.11.4），此理想与 $S'$ 的选择无关。这样得到的 $S$ 的理想就是 Zariski [11] 所称的“完备”理想。

\underline{定理} 2.13. \underline{在上述假设下，有}

\begin{equation*}
\chi(\underline{O}_S/I(\nu))
=\sum_{t\in S^*}\frac{\nu(t)(\nu(t)+1)}2[k(t):k],
\end{equation*}
\underline{以及}

\begin{equation*}
R^if_*\underline{O}_{S'}(-\nu^*)=0
\quad\text{对 }i>0.
\end{equation*}

用旧时的法语，这可简述为：“要求一条平面曲线以 $n$ 重经过一个实在的或无限邻近的点 $t$，便对其方程的系数施加 $n(n+1)/2$ 个一般相互独立的条件。”

先证明下述引理。

\underline{引理} 2.14. \underline{设 $p:S_1\longrightarrow S$ 是 $S$ 沿 $\{s\}$ 的爆破；设 $T$ 是 $S_1$ 的闭子概形，其理想为 $q$；设 $n$ 为整数；设 $J=\underline{O}_{S_1}(-s^*)$ 是定义除子 $s^*$ 的理想；置 $Y=\operatorname{Spec}(\underline{O}_{S_1}/J^nq)$；设 $D$ 是 $S$ 上的正则除子，其在 $S_1$ 上的严格变换为 $D'$。假设 $T\cap s^*$ 和 $T\cap D'$ 均有限。}

(i)\quad \underline{若 $\operatorname{length}_{k(s)}(T\cap s^*)\leq n+1$，则}

\begin{equation*}
R^ip_*(qJ^n)=0\quad\text{对 }i>0,
\end{equation*}
\underline{且若 $T$ 是 Artin 概形，则}

\begin{equation*}
\operatorname{length}_{k(s)}
\bigl(\underline{O}_S/p_*(qJ^n)\bigr)
=\frac{n(n+1)}2+\operatorname{length}_{k(s)}(T).
\end{equation*}

%%%%%%%%%% source scan idx 30 | folio 23 %%%%%%%%%%

(ii)\quad \underline{若 $\operatorname{length}_{k(s)}(T\cap s^*)\leq n$，则}

\begin{equation*}
\operatorname{length}_{k(s)}
\Bigl(D\cap
\operatorname{Spec}\bigl(\underline{O}_S/p_*(qJ^n)\bigr)\Bigr)
=n+\operatorname{length}_{k(s)}(D'\cap T) \quad .
\end{equation*}

考虑 $S_1$ 上的正合列

\begin{equation}\tag{2.14.1}
0\longrightarrow J^nq
\longrightarrow J^n
\xrightarrow{\,u\,}J^n\otimes\underline{O}_T
\longrightarrow0 \quad .
\end{equation}
由（2.11.4），要证明 $i>0$ 时 $R^ip_*(J^nq)=0$，只须证明 $p_*(u)$ 为满射。根据形式函数定理，只须证明各映射

\begin{equation*}
u_k:p_*(J^n\otimes\underline{O}/J^k)
\longrightarrow
p_*(J^n\otimes\underline{O}_T\otimes\underline{O}/J^k)
\end{equation*}
均为满射。

用 $\operatorname{Gr}_J$ 表示 $J$-进滤过的相伴分次对象。模 $J^m/J^{m+1}$ 同构于 $\underline{O}_{s^*}(m)$；因此 $R^1p_*(J^m/J^{m+1})=0$，且映射 \[
p_*(J^n/J^{n+k+1})\longrightarrow p_*(J^n/J^{n+k})
\] 为满射。所以，要对 $k$ 归纳证明各映射 $u_k$ 为满射，只须证明

\begin{equation*}
v_k:p_*\bigl(J^n\otimes\operatorname{Gr}_J^k(\underline{O})\bigr)
\longrightarrow
p_*\bigl(J^n\otimes\operatorname{Gr}_J^k(\underline{O}_T)\bigr)
\qquad(k\geq0).
\end{equation*}
的满射性。模 $\operatorname{Gr}_J^k(\underline{O}_T)$ 是 $J^k/J^{k+1}\otimes\underline{O}_T$ 的商，故其在 $k(s)$ 上的长度至多为 $n+1$。设 $N$ 为下列满射的核：

\begin{equation*}
v'_k:J^n\otimes\operatorname{Gr}_J^k(\underline{O})
\simeq\underline{O}_{s^*}(n+k)
\longrightarrow
J^n\otimes\operatorname{Gr}_J^k(\underline{O}_T)
\end{equation*}
它因而是射影直线 $s^*$ 上次数 $\geq-1$ 的可逆层。因此 $R^1p_*N=0$，且 $v_k=p_*(v'_k)$ 为满射。

若 $T$ 是 Artin 概形，（2.14.1）便给出

\begin{equation*}
\operatorname{length}_{k(s)}
\bigl(p_*J^n/p_*(J^nq)\bigr)
=\operatorname{length}_{k(s)}(\underline{O}_T),
\end{equation*}

%%%%%%%%%% source scan idx 31 | folio 24 %%%%%%%%%%

再由（2.11.4）及（2.11.6）即得 2.14(i)。

现在证明 (ii)。由于 $D$ 正则，其严格变换 $D'$ 同构于 $D$，且与 $s^*$ 相交的长度为一。因此假设 $\operatorname{length}_{k(s)}(T\cap s^*)\leq n$ 蕴含
\[
\operatorname{length}_{k(s)}\bigl((D'\cup T)\cap s^*\bigr)\leq n+1,
\]
所以 $D'\cup T$ 满足 (i) 的假设。

投影 $p$ 诱导 $D$ 与 $D'$ 之间的同构。在此同构下，$\underline{O}_D$ 中定义子概形
\[
D\cap\operatorname{Spec}\bigl(\underline{O}_S/p_*(qJ^n)\bigr)
\]
的理想，对应于 $D'$ 中由 $qJ^n$ 在 $S_1$ 上的整体截面生成的理想。此理想包含于 $D'$ 中定义 $D'\cap Y$ 的理想。事实上二者相等，因为 $\underline{O}_{D'\cup Y}$ 的每个截面均可延拓为 $\underline{O}_{S_1}$ 的截面：若 $I_{D'\cup Y}$ 是定义 $D'\cup Y$ 的理想，延拓的障碍位于 $R^1p_*(I_{D'\cup Y})$ 中。有

\begin{equation*}
\begin{aligned}
I_{D'\cup Y}
&=\underline{O}(-D')\cap qJ^n\\
&=(\underline{O}(-D')\cap J^n)\cap qJ^n\\
&=\underline{O}(-D')J^n\cap qJ^n\\
&=(\underline{O}(-D')\cap q)J^n,
\end{aligned}
\end{equation*}
由应用于 $D'\cup T$ 的 (i)，不存在障碍。

因此

\begin{equation}\tag{2.14.2}
\operatorname{length}_{k(s)}
\Bigl(D\cap
\operatorname{Spec}\bigl(\underline{O}_S/p_*(qJ^n)\bigr)\Bigr)
=\operatorname{length}_{k(s)}(D'\cap Y) \quad .
\end{equation}
同构

\begin{equation*}
J^n\otimes
\bigl(\underline{O}_{D'}/q\underline{O}_{D'}\bigr)
\xrightarrow{\;\sim\;}
\underline{O}_{D'}J^n/
\underline{O}_{D'}qJ^n
\end{equation*}
给出正合列

\begin{equation*}
0\longrightarrow
J^n\otimes\underline{O}_{D'\cap T}
\longrightarrow\underline{O}_{D'\cap Y}
\longrightarrow\underline{O}_{D'\cap ns^*}
\longrightarrow0,
\end{equation*}
从而

%%%%%%%%%% source scan idx 32 | folio 25 %%%%%%%%%%

\begin{equation}\tag{2.14.3}
\operatorname{length}_{k(s)}(D'\cap Y)
=n(D',s^*)+\operatorname{length}_{k(s)}(D'\cap T).
\end{equation}
由于 $(D',s^*)=1$，由（2.14.2）和（2.14.3）即得断言 2.14(ii)。

现在对 $\nu$ 的支集中元素的个数归纳，证明 2.13 的断言及下述断言：

\begin{equation}\tag{2.14.4}
\text{若 $D$ 是 $S$ 上的正则除子，则}\qquad
\chi\Bigl(D\cap\operatorname{Spec}(\underline{O}_S/I(\nu))\Bigr)
=\sum_{\substack{t\in S^*\\t\ \mathrm{on}\ D}}
\nu(t)[k(t):k].
\end{equation}
若 $\nu=0$，则 $\nu^*=0$、$I(\nu)=\underline{O}_S$，且 2.13 和（2.14.4）显然成立（2.13 的第二个断言可在 $n=0$ 时反复应用（2.11.4）得到）。否则 $\nu(s)>0$，且存在 $S$-态射 $g:S'\longrightarrow S_1$。

对 $S_1$ 的每个闭点 $a$，设 $S_{1a}$ 是 $S_1$ 在 $a$ 处的局部化，$\nu_a$ 是 $\nu$ 在 $S_{1a}^*$ 上的限制，令 $S'_a=S'\times_{S_1}S_{1a}$、$I(\nu_a)=g_*\underline{O}_{S'_a}(-\nu_a^*)$，并置
\[
T_a=\operatorname{Spec}(\underline{O}_{S_1}/I(\nu_a)),\qquad
T=\coprod_aT_a,
\]
其中的并有限且不交。设 $q$ 是 $T$ 的理想。沿用 2.14 的记号，有

\begin{equation}\tag{2.14.5}
g_*\underline{O}_{S'}(-\nu^*)=qJ^{\nu(s)}.
\end{equation}
归纳假设可应用于 $(S_{1a},S'_a,\nu_a)$。故

\begin{equation}\tag{2.14.6}
R^ig_*\underline{O}_{S'}(-\nu^*)=0\quad\text{对 }i>0,
\end{equation}
\begin{equation}\tag{2.14.7}
\chi(T)=
\sum_{t\in S^*\setminus\{s\}}
\frac{\nu(t)(\nu(t)+1)}2[k(t):k],
\end{equation}
又由应用于 $S_1$ 上除子 $s^*$ 的（2.14.4）及（2.12.1），

\begin{equation*}
\operatorname{length}_{k(s)}(T\cap s^*)
=\sum_{\substack{t\ \mathrm{proximate}\\\mathrm{to}\ s}}
\nu(t)[k(t):k(s)]
\leq\nu(s).
\end{equation*}

%%%%%%%%%% source scan idx 33 | folio 26 %%%%%%%%%%

由此不等式可以应用 2.14。从（2.14.5）得到

\begin{equation*}
R^ip_*g_*\underline{O}_{S'}(-\nu^*)=0
\quad\text{对 }i>0
\end{equation*}
以及
\begin{equation*}
\operatorname{length}_{k(s)}
\bigl(\underline{O}_S/p_*g_*\underline{O}_{S'}(-\nu^*)\bigr)
=\frac{\nu(s)(\nu(s)+1)}2+\operatorname{length}_{k(s)}(T).
\end{equation*}
因此 2.13 的断言由（2.14.6）及（2.14.7）得到。最后，2.14(ii) 给出

\begin{equation*}
\operatorname{length}_{k(s)}
\Bigl(D\cap\operatorname{Spec}(\underline{O}_S/I(\nu))\Bigr)
=\nu(s)+\operatorname{length}_{k(s)}(D'\cap T).
\end{equation*}
将（2.14.4）应用于 $D'$，得到

\begin{equation*}
\begin{aligned}
\chi\Bigl(D\cap\operatorname{Spec}(\underline{O}_S/I(\nu))\Bigr)
&=\nu(s)[k(s):k]
+\sum_{\substack{t\ \mathrm{on}\ D\\t\ne s}}
\nu(t)[k(t):k]\\
&=\sum_{\substack{t\in S^*\\t\ \mathrm{on}\ D}}
\nu(t)[k(t):k].
\end{aligned}
\end{equation*}
至此完成 2.13 的证明。